\documentclass[12pt]{article}                         
\pagestyle{plain} %-- Document Set Up Command

%-- Document Setup & Layout
\usepackage{fancyhdr}
\usepackage{comment}
\usepackage{multicol}
\usepackage[
  letterpaper,
  left=0.8in,
  right=0.8in,
  textheight=9.5in,
  bmargin=0.75in  % Adjust this value to push the footer down
]{geometry}

%-- Color & Text Formatting
\usepackage[svgnames]{xcolor} 
\usepackage{textcomp}
\usepackage{ulem}
\usepackage{graphicx}

%-- Math Packages
\usepackage{amsmath}     
\usepackage{amsfonts}
\usepackage{amstext}
\usepackage{amssymb}
\usepackage{latexsym}
\usepackage{bm}

%-- Graphics and Figures
\usepackage{graphicx}    
\usepackage{tikz}
\usepackage{pgfplots}
\usepackage{circledtext}
\usepackage{float}
\usepackage{caption}
\usepackage{subcaption}
\usetikzlibrary{decorations.markings, arrows.meta}
%-- Tables
\usepackage{array}      
\usepackage{tabularx}

%-- Lists
\usepackage{enumitem}    
%\usepackage{enumerate}
\usepackage{tasks}

%-- TikZ Libraries
\usetikzlibrary{shapes.geometric} 
\pgfplotsset{compat=1.18}
\tikzset{
    phase arrow/.style={
        postaction={decorate,
            decoration={
                markings,
                mark=at position 0.1 with {\arrow{Stealth}},
                mark=at position 0.3 with {\arrow{Stealth}},
                mark=at position 0.7 with {\arrow{Stealth}},
                mark=at position 0.9 with {\arrow{Stealth}}
            }
        }
    },
    phase arrow left/.style={
        postaction={decorate,
            decoration={
                markings,
                mark=at position 0.1 with {\arrow{Stealth[reversed]}},
                mark=at position 0.3 with {\arrow{Stealth[reversed]}},
                mark=at position 0.7 with {\arrow{Stealth[reversed]}},
                mark=at position 0.9 with {\arrow{Stealth[reversed]}}
            }
        }
    }
}
%-- Header/Footer
\pagestyle{fancy}
\fancyhf{} % Clear all header and footer fields
\fancyhead[L]{Jack Lingbeck} % Left header with name
\fancyhead[R]{May 13th, 2026} % Right header with date
\renewcommand{\headrulewidth}{0.4pt} % Horizontal line below the header
\rfoot{Page \thepage}
\lfoot{Math-675, Homework 5}

\begin{document}
\large
\begin{center}
\fbox{\begin{tabular}{c}
\bf Math 675 Spring 2026: \\
\bf Homework 4 (LATE)\\
\bf Due Tuesday, April 6th 2026\\
\end{tabular}}
\end{center}

\section*{Question 1:}
Consider the matrix $A = \begin{pmatrix} -5 & 1 \\ 4 & -2 \end{pmatrix}$ \\
The eigenvalues for A are
$\lambda_1 = -1$ and
$\lambda 2 = -6$ with corresponding eigenvectors 
$\mathbf{v}_{\lambda_1} = (1, 4)$ and 
$\mathbf{v}_{\lambda_2} = (-1, 1)$.
Use this and the method of undetermined coefficients to find the general solution to the inhomogeneous system:
\[\mathbf{x}'(t) = \begin{pmatrix} -5 & 1 \\ 4 & -2 \end{pmatrix} x(t) + e^{2t}\begin{pmatrix} 6 \\ -1 \end{pmatrix}\]

\section*{Question 2:}
Consider the matrix $A = \begin{pmatrix} 2 & -1 \\ 3 & 2 \end{pmatrix}$ \\
The eigenvalues for A are
$\lambda_1 = -1$ and
$\lambda_2 = 1$ with corresponding eigenvectors 
$\mathbf{v}_{\lambda_1} = (1, 3)$ and 
$\mathbf{v}_{\lambda_2} = (1, 1)$.

\begin{enumerate}  
    \item \textbf{Calculate} $\boldsymbol{\exp(tA)}$ \\
    Given the definition: $\exp(tA) = \Psi(t) \Psi^{-1}(0)$, where $\Psi(t)$ is the fundamental matrix made up of the $x$ and $y$ solution vectors from $\mathbf{x}(t)$. \\
    First, build the general solution: 
    \[ \mathbf{x}(t) = c_1 e^{-t} \begin{pmatrix} 1 \\ 3 \end{pmatrix} + c_2 e^{t} \begin{pmatrix} 1 \\ 1 \end{pmatrix} \]
    \[\Psi(t) = \begin{pmatrix} e^{-t} & e^{t} \\ 3e^{-t} & e^{t}\end{pmatrix}\]
    \[\Psi^{-1}(0) = (\Psi(0))^{-1} = \begin{pmatrix} 1 & 1 \\ 3 & 1 \end{pmatrix}^{-1} = \frac{1}{2} \begin{pmatrix} -1 & 1 \\ 3 & -1 \end{pmatrix}\]
    \[ \exp(tA) = \Psi(t)\Psi^{-1}(0) = \begin{pmatrix} e^{-t} & e^{t} \\ 3e^{-t} & e^{t}\end{pmatrix} \frac{1}{2} \begin{pmatrix} -1 & 1 \\ 3 & -1 \end{pmatrix}\]
    \[ \exp(tA) = \frac{1}{2} \begin{pmatrix} -e^{-t} + 3e^{t} & e^{-t} - e^{t} \\ -3e^{-t}+3e^{t} & 3e^{-t} - e^{t}\end{pmatrix}\]
    Now, optionally, we could convert this into a form involving hyperbolic trig:
    \[ \exp(tA) = \begin{pmatrix} \sinh(t) & -\sinh(t) \\ 3\sinh(t) & \cosh(t) - 2\sinh(t) \end{pmatrix}\]
    \item \textbf{Solve the IVP} \\
    Use the result from part (1) to solve:
    \[ \mathbf{x}'(t) = \begin{pmatrix} 2 & -1 \\ 3 & -2 \end{pmatrix} \mathbf{x}(t), \quad \mathbf{x}(0) = \begin{pmatrix} 1 \\ 2 \end{pmatrix} \]
    Defintion: Given a system $\mathbf{x}'(t) = Ax(t)$ the solution is $\mathbf{x}(t)=e^{tA}x(0)$
    \[\mathbf{x}(t) = e^{tA}x(0) = \frac{1}{2} \begin{pmatrix} -e^{-t} + 3e^{t} & e^{-t} - e^{t} \\ -3e^{-t}+3e^{t} & 3e^{-t} - e^{t}\end{pmatrix} \begin{pmatrix} 1 \\ 2 \end{pmatrix}\]
    \[\mathbf{x}(t) = e^{tA}x(0) = \frac{1}{2} \begin{pmatrix} -e^{-t} + 3e^{t} + 2e^{-t} - 2e^{t} \\ -3e^{-t} + 3e^{t} + 6e^{-t} - 2e^{t} \end{pmatrix}\]
    \[\mathbf{x}(t) = e^{tA}x(0) = \frac{1}{2} \begin{pmatrix} e^{-t} + e^{t} \\ 3e^{-t} + e^{t} \end{pmatrix}\]
    \[\mathbf{x}(t) = \frac{1}{2}e^{-t}\begin{pmatrix} 1 \\ 3 \end{pmatrix} + \frac{1}{2} e^{t} \begin{pmatrix} 1 \\ 1 \end{pmatrix}\]
    Ta-Da! This is the same IVP Solution as if we solved this the regular way!
\end{enumerate}

\section*{Question 3:}
The following questions refer to the proof of existence and uniqueness covered in chapter 2.3 of the guided notes.
\begin{enumerate}
    \item[a)] \textbf{Why do we choose a small time interval?} \\
    In the proof we chose $\alpha < min\left\{\frac{\rho}{M},\frac{1}{B}\right\}$
    \begin{enumerate}
        \item Explain in words why we need the condition $\alpha < \frac{\rho}{M}$
        \item Explain in words why we need the condition $\alpha < \frac{1}{B}$
        \item What could go wrong if $\alpha$ was too large?
    \end{enumerate}

    \item[b)] \textbf{Staying inside the closed ball} \\
    Suppose we have $x_{k+1}(t) = x_0 + \int_{0}^{t} F(x_k(x))ds$
    \begin{enumerate}
        \item Why is it important that all iterates stay inside $\mathcal{O}_{\rho}$
        \item What assumptions about $F$ would fail if an iterate left $\mathcal{O}_{\rho}$?
        \item Explain how the bound $\Vert F(w) \Vert \leq M$ is used in this step
    \end{enumerate}

    \item[c)] \textbf{Role of Lipschitz Condition} \\
    The proof uses $\Vert F(x) - F(y) \Vert \leq B \Vert x - y \Vert$
    \begin{enumerate}
        \item Where in the proof is this condition used?
        \item Why is continuity of $F$ alone not enough for the argument?
        \item In your own words, what does the Lipschitz condition guarentee?
    \end{enumerate}

    \item[d)] \textbf{Why does a geometric series appear?} \\
    We obtained an estimate $\Vert x_{k+1}(t) - x_k(t) \Vert \leq (\alpha B)^k L$
    \begin{enumerate}
        \item Why is it important that $\alpha B < 1$?
        \item What would happen if $\alpha B \geq 1$
        \item How does this estimate lead to convergence of $\{x_k\}$
    \end{enumerate}

    \item[e)] \textbf{Conceptual comparison} \\
    Compare:
    \begin{itemize}
        \item Solving an algebric equation $x = G(x)$ via fixed point iteration
        \item Solving the integral equation $x(t) = x_0 + \int_{0}^{t} F(x(s))ds$
    \end{itemize}
    \begin{enumerate}
        \item What plays the role of the function being iterated?
        \item What plays the role of the fixed point?
        \item Why is this analogy useful?
    \end{enumerate}
\end{enumerate}

\end{document}

